\section{Introduction}
\label{sec:intro}

% Paragraph plan (one idea per paragraph):
% 1. Magnets as cost/size driver of a fusion device: field on coil limits
%    B0, hence fusion power and machine size. HTS opening new design space
%    (compact high-field tokamaks, VNS, ...). \todo{refs}
% 2. Where magnet design sits today: systems codes (PROCESS-like, with
%    simplified magnet models) vs. detailed engineering design (FEM, coil
%    by coil, months). Gap: a fast but physically grounded tool that maps
%    a configuration into a buildable magnet. \todo{refs to systems codes
%    and to magnet pre-design tools}
% 3. Key idea of MADE: the magnetic configuration is the INPUT. The user
%    fixes what the plasma physicist asks (B0, R0, A, ripple, flux swing,
%    PF currents, mirror ratio...) and MADE returns the magnets
%    (conductor, winding pack, structure) that deliver it.
% 4. Second idea: explore, do not optimise. All feasible design points
%    are kept -> trade-off maps (radial build vs current vs Jeng vs
%    stress vs hot spot), robustness to assumptions.
% 5. Third idea: one set of modules, many devices (tokamak TF/CS/PF,
%    mirror, stellarator in progress).
% 6. History: MADE was born for the DTT CS [Giannini2020DTT, Giannini2021DTT]
%    and extended in EUROfusion DEMO: macro-blocks I (TF) and II (CS)
%    [Giannini2023MADE], macro-block III (TF shape optimisation + PF,
%    ITER case) [Giannini2024MADEIII]. Companion analytical work:
%    [Portone2023HTS].
%    What is new in 2.0:
%    - one modular code base (io/physics/search/postprocess) per coil
%      family, Excel input, no machine data in the code;
%    - TF: peak field per grade from a discrete 2D field model calibrated
%      at scan start and validated against ANSYS (replaces the linear
%      smeared gradient + constant correction of v1);
%    - TF: FE-calibrated mechanical surrogates for jacket and case;
%    - CS/PF: axisymmetric FE verification of the chosen design inside
%      the code, scenario forces, FCGR fatigue as sizing option;
%    - 3D module: bending-free shape iteration with Biot-Savart on the
%      discrete coil set, and a global beam+shell model of the whole TF
%      system verified against ANSYS;
%    - mirrors with the same modules; path to stellarators.
% 7. Paper outline.

\todo{Write from the paragraph plan above.}

The rest of the paper is organised as follows.
Section~\ref{sec:philosophy} describes the design philosophy and the
architecture of \MADE{}. Section~\ref{sec:conductor} presents the
coil-agnostic conductor module. Sections~\ref{sec:tf}
and~\ref{sec:cspf} describe the TF and the CS/PF solvers, and
section~\ref{sec:mirror} shows how they are reused for mirror machines.
Section~\ref{sec:stellarator} introduces the global beam--shell model and
the path towards stellarator coils. Verification and applications are
given in sections~\ref{sec:verification} and~\ref{sec:applications}.
